\documentclass[10pt,a4paper]{amsart}
\usepackage[margin=19mm]{geometry}
\usepackage{amsmath,amssymb,mathtools}
\usepackage[scr=rsfs,cal=euler]{mathalfa}
\usepackage{xfrac}
\usepackage[unicode,hidelinks,bookmarks=false]{hyperref}
\newcommand{\ie}{i.e.\ }
\newcommand{\cf}{cf.\ }
\newcommand{\resp}{respectively\ }
\newcommand{\Subsection}{\subsection}
\providecommand{\nobreakdash}{\nobreak\hbox{-}}
\setlength{\emergencystretch}{2em}
\multlinegap0pt
\usepackage{dsfont}
\usepackage{amssymb}
\usepackage{mathtools}
\usepackage{xcolor}
\usepackage{tikz}
\usepackage[shortlabels]{enumitem}
\newtheorem{corollary}{Corollary}
\newcommand{\NN}{{\mathcal N}}
\newcommand{\e}{\epsilon}
\newcommand{\mf}{\mathfrak}
\title{Partition regularity in imaginary quadratic rings of integers}
\author{Sebasti\'an Donoso, Andreu Ferr\'e Moragues, Andreas Koutsogiannis,, Wenbo Sun}
\date{Source: arXiv:2603.20497v2 (CC BY 4.0)}
\begin{document}
\maketitle

\noindent\emph{Source and licence.} Partition regularity in imaginary quadratic rings of integers. Version arXiv:2603.20497v2 (CC BY 4.0), distributed by arXiv under the Creative Commons Attribution 4.0 International licence, http://creativecommons.org/licenses/by/4.0/, as recorded in the arXiv licence metadata for this exact version and shown on the abstract page https://arxiv.org/abs/2603.20497v2. Full licence notice: \texttt{licenses/arxiv-2603.20497-CC-BY-4.0.txt}.\par\smallskip

\noindent\textit{Editorial scope.} Counted unit: the article's corollary x (source0.tex lines 1587--1593). The article's complete original proof of it is delivered with it (source0.tex lines 2565--2657, one \texttt{proof} environment of 93 lines, which the article places away from the statement). No mathematics is written, repaired or re-derived: the statement and the proof are the article's own lines, in the article's own notation, and no retained line is substituted or re-flowed.  Retained uncounted as the source-local context the proof needs: lines 2548--2550; lines 2552--2554.  Nothing was omitted from any retained span, so the package carries no omission record.  No retained line contains a citation, so the wrapper carries no bibliography and no bibliography entry.  No retained line contains a figure environment, an image-inclusion command or drawing code, so no image is delivered and the package declares no asset dependency.  Editorial formatting only, in addition: an AMS \texttt{amsart} preamble that restates the article's own declarations and macros used by the retained lines, namely the environments \texttt{corollary} on separate counters, together with the standard packages those lines need.  The retained environments print in this document as Corollary 1 in order of appearance, whereas the article prints the same items with the numbering of its own full document.  The complete native source of the article as distributed in the arXiv e-print for this exact version is bundled unchanged as \texttt{sources/196849/source0.tex}.
 \begin{equation}\label{1}
     \sum_{1\leq\NN(\mf u)\leq x}g(\mf u)=\frac{x^{1+i\alpha}}{1+i\alpha}CL(\log  x)+o(x), \text{ as }x\to\infty
 \end{equation}

  \begin{equation}\label{2}
     G(s)=\frac{C}{s-(1+i\alpha)}L\left(\frac{1}{\sigma-1}\right)+o\left(\frac{1}{\sigma-1}\right), \text{ as }\sigma\to 1
 \end{equation}

\begin{corollary}\label{lem: bounds for number of ideals of norm at most x}
Let $K$ be a number field of degree $D$. Then, there exists a universal constant $\gamma_K>0$ and some $0 \leq \eta<1$ %and some $0<\delta<1$ 
such that for any ideal class $C$ of $K$, we have that %the following holds:
\[
\frac{1}{x}\sum_{\mf u\in C\colon \NN(\mf u) \leq x}1 = \gamma_K+o(x^{\eta}).
\]
\end{corollary}

  \begin{proof}[Proof of \eqref{1} implies \eqref{2}] 
 % Let
 % $$H(x):=\sum_{1\leq \NN(\mf u)\leq x}g(\mf u).$$
 Note that $H(1-)=0$ and 
    $\lim_{N\to\infty}\left| \frac{H(x)}{sx^{s}}\right|\leq \lim_{N\to\infty}\frac{O(1)}{sx^{\sigma-1}}=0$ if $\sigma>1$ by Corollary~\ref{lem: bounds for number of ideals of norm at most x}.    
    Using the properties of the Riemann-Stieltjes integral,
   %(write $u=u_{1}+iu_{2}$)
    \begin{equation}\label{early}
    \begin{split}
    &\qquad \frac{G(s)}{s}=\frac{1}{s}\lim_{R\to\infty}\int_{1_-}^R\frac{1}{x^{s}}\,dH(x)=-\lim_{R\to\infty}\int_{1}^R H(x)\,d\left(\frac{1}{sx^{s}}\right)
    \\&\qquad\;\;\;\quad\quad=\lim_{R\to\infty}\int_1^R\frac{1}{x^{s+1}}H(x)\,dx = \int_1^{\infty} \frac{1}{x^{s+1}}H(x)\,dx.
    \end{split}
 \end{equation}
 %  Indeed, it is clear that $H(1-)=0$ by construction, and $\left| \frac{H(x)}{su^{s}}\right|\leq C_K\frac{R}{sR^{\sigma}}$, for some constant $C_K$ depending on the number field $\mathcal{O}_K$   provided that $x\geq R$ which, in turn, goes to $0$ as $R \to \infty$ with our choice of $s$.
   
Let $\varepsilon>0$. By (\ref{1}), there exists $x_{0}\geq 1$ such that if $x \geq x_{0}$, then we have that 
 $$\vert H(x)-\lambda x^{1+i\alpha}L(\log x)\vert<\varepsilon x,$$
 where $\lambda=C/(1+i\alpha)$. Then
     \begin{equation}\nonumber
    \begin{split}
    &\qquad \left| \int_{x \geq x_{0}}\frac{1}{x^{s+1}}H(x)\,dx-\int_{x\geq x_{0}}\frac{1}{x^{s+1}}\lambda x^{1+i\alpha}L(\log x) \ dx\right|
   \leq %\int_{\max\{u_{1},u_{2}\}\geq x_{0}} \e \NN(x)^{-\sigma}\,d(\NN(x))
    \varepsilon\int_{x\geq 1} x^{-\sigma}\,dx = \frac{\varepsilon}{\sigma-1}.
    %=\varepsilon\frac{1}{\sigma-1}\left(-\int_{\NN(u)\geq 1} \,d(\NN(u)^{1-\sigma})\right)
    %\leq \frac{\varepsilon C_0}{\sigma-1},
    \end{split}
 \end{equation} 
%for some universal constant $C_{0}$. 
  Therefore, we have that 
  \begin{equation}\nonumber
    \begin{split}
    \frac{G(s)}{s}=\lambda \int_{x\geq 1} x^{-s+i\alpha}L(\log  x)\,dx+o\left(\frac{1}{\sigma-1}\right).
    \end{split}
 \end{equation}
 Fix $\Delta>1$. Let $y_{1}=\exp(\Delta^{-1}(\sigma-1)^{-1})$ and  $y_{2}=\exp(\Delta(\sigma-1)^{-1})$. Then since $L$ is slowly oscillating, we have that
 $$L(\log  x)=L\left(\frac{1}{\sigma-1}\right)+o(1),$$
 uniformly for all $y_{1}\leq  x\leq y_{2}$. Therefore, 
   \begin{equation}\nonumber
    \begin{split}
   &\qquad \int_{y_{1}\leq x\leq y_{2}}x^{-s+i\alpha}L(\log x)\,dx
   =\int_{ y_{1}\leq x\leq y_{2}}x^{-s+i\alpha}L\left(\frac{1}{\sigma-1}\right)\,dx+o\left(\frac{1}{\sigma-1}\right).
  % \\&=\int_{ y_{1}\leq \NN(x)\leq y_{2}}\frac{1}{\NN(x)^{s+1}}\omega(x)^{1+i\alpha}L\left(\frac{1}{\sigma-1}\right)\,d(\NN(x))+,
    \end{split}
 \end{equation}
Therefore, given that $|L|=1$, using crude upper bounds, as well as the computation above, we have
  \begin{eqnarray*}
      \left|G(s)-\frac{s\lambda}{s-(1+i\alpha)} L\left(\frac{1}{\sigma-1}\right) \right|
   & = & \vert s\vert\cdot\left|\frac{G(s)}{s}-\lambda L\left(\frac{1}{\sigma-1}\right) \int_{x\geq 1}x^{-s+i\alpha}\,dx\right|
   \\ &\leq & 2\vert s\lambda\vert\int_{x<y_{1} \text{ or } x>y_{2}} x^{-\sigma}\,dx+o\left(\frac{1}{\sigma-1}\right)\\
 & \leq & \frac{1}{\sigma-1}(2\vert s\lambda\vert(\Delta^{-1}+e^{-\Delta})+o(1)).
  \end{eqnarray*}
Since $\vert L\vert=1$ and 
\[
 \frac{s\lambda}{s-(1+i\alpha)} = \lambda+\frac{C}{s-(1+i\alpha)},
 \]
 we have that 
  \begin{equation}\nonumber
    \begin{split}
   \left|G(s)-\frac{C}{s-(1+i\alpha)} L\left(\frac{1}{\sigma-1}\right) \right|
  \leq \frac{1}{\sigma-1}(2\vert s\lambda\vert(\Delta^{-1}+e^{-\Delta})+o(1)).
       \end{split}
 \end{equation}
 
   %\\&=\frac{1}{\sigma-1}(2\vert\lambda\vert((1-e^{-\Delta^{-1}})^{2}+(2e^{-0.5\Delta}))+o(1))
 %  \leq \frac{1}{\sigma-1}(2\vert\lambda\vert(\Delta^{-2}+(2e^{-0.5\Delta}))+o(1)).
% {\color{red} Here we view $\int \NN(u)^{-\sigma}\,d(\NN(u))$ as $-\frac{1}{\sigma-1}\int d(\NN(u)^{1-\sigma})$ and do integration by parts. For the part $\NN(u)<y_{1}$, we pick the end points as $1+i$, $1+i\sqrt{y_{1}-1}$, $\sqrt{y_{1}-1}+i$ and $\sqrt{y_{1}-1}+i\sqrt{y_{1}-1}$. For the part $\NN(u)>y_{2}$, we pick the end points of the two infinite boxes as $\sqrt{y_{2}}/2$ and  $\sqrt{y_{2}}i/2$.
% }
% {\color{blue} Careful here with the norm of $1+i$ being equal to $\sqrt{2}$. I think this changes some of the bounds by constants in the third and fourth lines, but after all, since we are going to send $\Delta$ to infinity, this will not matter. I think the result we care about holds modulo constants.}
 Let $\sigma$ approach $1$ from the right first, and then let $\Delta$ grow very large, we get (\ref{2}).
 %
 \iffalse
 
 . It follows from the computation above that
 \[
  G(s)=\frac{\lambda}{s-(1+i\alpha)} L\left(\frac{1}{\sigma-1}\right) +o\left(\frac{1}{\sigma-1}\right)
 \]
 so we complete this part of the proof by using the following two facts. First, that $|L|=1$, and second, that
 \[
 \frac{s\lambda}{s-(1+i\alpha)} = \lambda+\frac{C}{s-(1+i\alpha)}
 \]
  % \begin{equation}\nonumber
  %  \begin{split}
  % &\qquad\int_{u_{1},u_{2}\geq 1}\frac{s}{\NN(u)^{s+1}}\omega(u)^{1+i\alpha}\,d(\NN(u))+o\left(\frac{1}{\sigma-1}\right)
  % \\&=\frac{s\lambda}{s-(1+i\alpha)} L\left(\frac{1}{\sigma-1}\right) \int_{u_{1},u_{2}\geq 1}\,d(-\NN(u)^{(1+i\alpha)-s})+o\left(\frac{1}{\sigma-1}\right)
     % \\&=-\lambda L(\frac{1}{\sigma-1}) \int_{u_{1},u_{2}\geq 1}\omega(u)^{1+i\alpha}\,d(\NN(u)^{-s})+o(\frac{1}{\sigma-1})
   %    \\&=\frac{s\lambda}{s-(1+i\alpha)}K(s-(1+i\alpha))  L\left(\frac{1}{\sigma-1}\right)+o\left(\frac{1}{\sigma-1}\right)
    %    \\&=(\lambda+\frac{C}{s-(1+i\alpha)})K(s-(1+i\alpha)) L\left(\frac{1}{\sigma-1}\right)+o\left(\frac{1}{\sigma-1}\right)
     %   \\&=\frac{C}{s-(1+i\alpha)}K(s-(1+i\alpha)) L\left(\frac{1}{\sigma-1}\right)+o\left(\frac{1}{\sigma-1}\right).
    %\end{split}
 %\end{equation}
 %{\color{red} take care of this constant $K(s-(1+i\alpha))=2^{s-(1+i\alpha)}$ in the statement of the theorem}
\fi
  \end{proof}

\end{document}
